% !TEX root = ../../main.tex
\section{From ZooKeeper to KRaft: Kafka's Metadata Evolution~\cite{kafka_kraft}}

% TODO: Explain why ZooKeeper was originally used, what problems it caused
%       (separate cluster to operate, scalability ceiling, split-brain risk),
%       and how KRaft (KIP-500) eliminates the dependency entirely.

\subsection{The ZooKeeper Era}

For its first decade, Kafka relied on \textbf{Apache ZooKeeper} to manage all
cluster metadata:

\begin{itemize}
  \item \textbf{Broker registration} — each broker registered itself in ZooKeeper on startup; clients discovered brokers via ZooKeeper znodes.
  \item \textbf{Controller election} — one broker was elected the \emph{controller} via a ZooKeeper ephemeral node; it was responsible for leader elections and partition reassignments.
  \item \textbf{Topic and partition metadata} — partition-to-broker assignments, ISR lists, and configuration overrides were stored in ZooKeeper.
\end{itemize}

\begin{figure}[H]
  \centering
  \fbox{\parbox{0.85\textwidth}{\centering\vspace{2.5cm}[Kafka + ZooKeeper Architecture]\vspace{2.5cm}}}
  \caption{Kafka cluster with ZooKeeper for metadata management (legacy mode).}
  \label{fig:kafka-zk}
\end{figure}

\textbf{Pain points of the ZooKeeper dependency:}
\begin{itemize}
  \item Operators had to maintain a \emph{separate} ZooKeeper ensemble (typically 3–5 nodes) alongside the Kafka cluster.
  \item ZooKeeper became a scalability bottleneck: clusters with millions of partitions overwhelmed the ZooKeeper write path.
  \item Controller failover required a full ZooKeeper session timeout, causing 20–30 second leadership gaps.
  \item Security and monitoring configurations were duplicated across two systems.
\end{itemize}

\subsection{The KRaft Era (Kafka without ZooKeeper)}

\textbf{KIP-500} (released preview in Kafka 2.8, production-ready in Kafka 3.3)
introduced \textbf{KRaft} — Kafka Raft — a built-in consensus protocol that
replaces ZooKeeper entirely.

\begin{figure}[H]
  \centering
  \fbox{\parbox{0.85\textwidth}{\centering\vspace{2.5cm}[Kafka KRaft Architecture]\vspace{2.5cm}}}
  \caption{Kafka cluster in KRaft mode: no external ZooKeeper required.}
  \label{fig:kafka-kraft}
\end{figure}

In KRaft mode:
\begin{itemize}
  \item A subset of brokers are designated \textbf{controllers}; they form a Raft quorum and store all metadata in an internal \textbf{metadata log} (a special Kafka topic: \texttt{\_\_cluster\_metadata}).
  \item All other brokers are \textbf{followers} that subscribe to the metadata log and maintain a local in-memory cache — no ZooKeeper znodes, no separate cluster.
  \item Controller failover is now a \textbf{Raft leader election}, completing in milliseconds instead of tens of seconds.
  \item The metadata log can handle millions of partitions without the write-path bottleneck seen with ZooKeeper.
\end{itemize}

\begin{table}[H]
  \centering
  \caption{ZooKeeper mode vs KRaft mode.}
  \label{tab:zk-vs-kraft}
  \begin{tabularx}{\textwidth}{lXX}
    \toprule
    \textbf{Dimension}        & \textbf{ZooKeeper Mode}               & \textbf{KRaft Mode} \\
    \midrule
    External dependency       & ZooKeeper ensemble (3–5 nodes)        & None \\
    Metadata storage          & ZooKeeper znodes                      & Internal \texttt{\_\_cluster\_metadata} topic \\
    Controller failover time  & 20–30 seconds (ZK session timeout)    & Milliseconds (Raft leader election) \\
    Partition scalability     & $\sim$200k partitions practical limit & Millions of partitions \\
    Operational complexity    & Two systems to configure \& monitor   & Single Kafka cluster \\
    Security model            & Separate ACLs on ZooKeeper            & Unified Kafka ACLs \\
    Production readiness      & Deprecated (Kafka 3.5+)               & Production-ready from Kafka 3.3 \\
    \bottomrule
  \end{tabularx}
\end{table}
